
\subsection{Clebsch-Gordan Coefficients}

De volgende eigenschappen kunnen worden aangetoond:

\textbf{Orthonormaliteit:}
\[
\sum_{i_1, i_2} \bar{C}_{i_1 i_2 i_3 s}^{\alpha \beta \gamma} \, C_{i_1 i_2 i_3' s'}^{\alpha \beta \gamma'} = \delta_{\gamma, \gamma'} \, \delta_{s, s'} \, \delta_{i_3, i_3'}
\]

\textbf{Compleetheid:}
\[
\sum_{\gamma, s, i_3} \bar{C}_{i_1 i_2 i_3 s}^{\alpha \beta \gamma} \, C_{i_1' i_2' i_3 s}^{\alpha \beta \gamma} = \delta_{i_1, i_1'} \, \delta_{i_2, i_2'}
\]

\begin{proof}
De Clebsch-Gordan coëfficiënten zijn per definitie de matrixelementen van de unitaire transformatie die de productbasis $\{\langle{i_1 i_2}\}$ van $V_\alpha \otimes V_\beta$ koppelt aan de irreducibele basis $\{\ket{\gamma s i_3}\}$, waarin de productrepresentatie blokdiagonaal ontbonden is:
\[
\ket{\gamma s i_3} = \sum_{i_1,i_2} C^{\alpha\beta\gamma}_{i_1 i_2 i_3 s} \ket{i_1 i_2},
\qquad
\ket{i_1 i_2} = \sum_{\gamma,s,i_3} \bar{C}^{\alpha\beta\gamma}_{i_1 i_2 i_3 s} \ket{\gamma s i_3}.
\]

Omdat beide bases orthonormaal zijn, is $C$ een \emph{unitaire} matrix:
\[
C^\dagger C = \mathds{1}, \qquad C C^\dagger = \mathds{1}.
\]

\textbf{Orthonormaliteit.}
Schrijf de eerste relatie ($C^\dagger C = \mathds{1}$) uit in indices. De rij-index van $C$ is $(i_1,i_2)$, de kolom-index is $(\gamma,s,i_3)$. Dus
\[
(C^\dagger C)_{(\gamma s i_3),(\gamma' s' i_3')}
= \sum_{i_1,i_2} \bar{C}^{\alpha\beta\gamma}_{i_1 i_2 i_3 s} \, C^{\alpha\beta\gamma'}_{i_1 i_2 i_3' s'}
= \delta_{(\gamma s i_3),(\gamma' s' i_3')}
= \delta_{\gamma\gamma'}\,\delta_{ss'}\,\delta_{i_3 i_3'}.
\]
We hebben de orthonormaliteitsrelatie. 

Equivalent: $\braket{\gamma s i_3 | \gamma' s' i_3'} = \delta_{\gamma\gamma'}\delta_{ss'}\delta_{i_3i_3'}$, geschreven in de productbasis.

\textbf{Compleetheid.}
Tweede relatie ($CC^\dagger = \mathds{1}$):
\[
(C C^\dagger)_{(i_1 i_2),(i_1' i_2')}
= \sum_{\gamma,s,i_3} C^{\alpha\beta\gamma}_{i_1 i_2 i_3 s} \, \bar{C}^{\alpha\beta\gamma}_{i_1' i_2' i_3 s}
= \delta_{(i_1 i_2),(i_1' i_2')}
= \delta_{i_1 i_1'}\,\delta_{i_2 i_2'}.
\]
Equivalent: $\braket{i_1 i_2 | i_1' i_2'} = \delta_{i_1i_1'}\delta_{i_2i_2'}$, uitgedrukt in de irreducibele basis.

De eigenschappen zijn dus de twee kanten van de unitariteit van de basistransformatie $C$ uitwerken.
\end{proof}



\subsubsection{Relaties met Representaties}
Een gerelateerde, vaak gebruikte identiteit is de volgende decompositie van de productrepresentatie, die ook uit unitariteit van $C$ volgt:
\[
D^\alpha_{i_1 i_1'}(g)\, D^\beta_{i_2 i_2'}(g)
= \sum_{\gamma,s,i_3,i_3'} C^{\alpha\beta\gamma}_{i_1 i_2 i_3 s}\, D^\gamma_{i_3 i_3'}(g)\, \bar{C}^{\alpha\beta\gamma}_{i_1' i_2' i_3' s}.
\]
Combineren met Schur levert
\[
\frac{1}{\#G}\sum_{g\in G} D^\alpha_{i_1 i_1'}(g)\, D^\beta_{i_2 i_2'}(g)\, \bar{D}^\gamma_{i_3 i_3'}(g)
= \frac{1}{d_\gamma}\sum_{s} C^{\alpha\beta\gamma}_{i_1 i_2 i_3 s}\, \bar{C}^{\alpha\beta\gamma}_{i_1' i_2' i_3' s}.
\]
Deze relatie is interessant. (Zoals Frank of Bram het zou zeggen)